\documentclass[11pt]{article}
\usepackage[margin=25mm]{geometry}
\usepackage{fontspec}
\setmainfont{Latin Modern Roman}
\usepackage{amsmath,amssymb,mathtools}
\usepackage{unicode-math}
\setmathfont{Latin Modern Math}
\usepackage{xurl,hyperref}
\hypersetup{colorlinks=true,urlcolor=blue}
\setcounter{secnumdepth}{0}
\setlength{\parindent}{0pt}
\setlength{\parskip}{5pt}
\setlength{\emergencystretch}{3em}
\newcommand{\sref}[2]{\hyperlink{source:#1:#2}{[S#2]}}
\newcommand{\cataloguescope}[1]{\paragraph{Recorded target.}#1}
\begin{document}
\section{Research Problems in Function Theory — Problem 6.61}\label{statement}

\subsection{Definitions and mathematical statement}

\paragraph{Shared notation.}$\mathbb N=\{1,2,\ldots\}$, and $\mathbb Z,\mathbb Q,\mathbb R,\mathbb C$ denote the integers, rationals, reals and complex numbers. Any other conventions are specified locally.


A Jordan domain is the interior of a simple closed curve. A rectifiable curve has finite arclength. Let $D_1,D_2$ be Jordan domains whose boundary curves $C_1,C_2$ have the same finite length. An isometric sewing identifies their boundary points by an orientation-reversing arclength-preserving homeomorphism. It is conformally admissible everywhere if the sewn topological sphere has a Riemann-surface structure compatible with the original complex structures on both interiors. Uniformize this surface conformally by the Riemann sphere; denote the image of the seam by $C$. Rectifiability of $C$ is measured in the round spherical metric (equivalently in a smooth planar coordinate avoiding $C$).
\paragraph{Statement recorded in the supplied source.}
For every everywhere conformally admissible isometric sewing as above, must the seam $C$ be rectifiable? \sref{statement}{2}
\subsection{Short English statement}
Does sewing two equal-length rectifiable Jordan boundaries conformally always give a rectifiable seam on the uniformized sphere?
\subsection{Sources}
\begin{description}
\item[\hypertarget{source:statement:1}{S1}] ULAM.AI and the UnsolvedMath Contributors, \emph{UnsolvedMath: A Curated Collection of Open Mathematics Problems}, catalogue entry. CC BY 4.0. \url{https://huggingface.co/datasets/ulamai/UnsolvedMath}.
\item[\hypertarget{source:statement:2}{S2}] W. K. Hayman and E. F. Lingham, \emph{Research Problems in Function Theory} (2018), Chapter 6, Problem 6.61 and its complete update; Chapter 6 initial notation. \url{https://arxiv.org/abs/1809.07200}.
\end{description}
\label{end:statement}
\end{document}
